% GENERATED FILE. Edit canonical lessons, then run npm run generate:books.
% canonical-source-hash: fnv1a64:1a733a9bcff06720
% canonical-lessons: LA-C135-balaena, LA-C135-gallus, LA-C135-labrum, LA-C135-mentum, LA-C135-dorsum

\chapter{Whale, Cockerel, Lip, Chin, Back}
\label{ch:la-whale-cockerel-lip-chin-back}
\hlchaptermodality{135}

\begin{chapteropening}
\textbf{By the end of this chapter:} \emph{I can say a whale, a cockerel, a lip, the chin and the back.}

Complete the last lesson of chapter 135: i can say a whale, a cockerel, a lip, the chin and the back.
\end{chapteropening}

\section[balaena, balaenae]{balaena, balaenae — a whale}

\label{lesson:LA-C135-balaena}



\begin{quote}
Before the new one: say the Latin for a monkey, then the Latin for a dolphin.
\end{quote}

\subsection*{You'll want to know: balaena, balaenae}
\textbf{balaena, balaenae} — \textquotedblleft{}a whale\textquotedblright{}.

\begin{grammarlens}[title={What you've built}]
One more word for this chapter.
\end{grammarlens}

\subsection*{Your turn}
\begin{itemize}
\raggedright
  \item \emph{Say it:} \emph{balaena}
  \item \emph{Say it:} \emph{balaena}, once more
  \item \emph{Say it:} say \emph{balaena}
  \item \emph{Recall:} say the Latin for a monkey, then the Latin for a dolphin, then say \emph{balaena} again
\end{itemize}

\begin{tcolorbox}[breakable,colback=teal!4,colframe=teal!35!black,title={Before you move on}]
What does balaena, balaenae mean? (\textquotedblleft{}A whale\textquotedblright{}.) Say it once more.
\end{tcolorbox}
\section[gallus, gallī]{gallus, gallī — a cockerel}

\label{lesson:LA-C135-gallus}



\begin{quote}
Before the new one: say the Latin for a dolphin, then the Latin for a whale.
\end{quote}

\subsection*{You'll want to know: gallus, gallī}
\textbf{gallus, gallī} — \textquotedblleft{}a cockerel\textquotedblright{}.

\begin{grammarlens}[title={What you've built}]
One more word for this chapter.
\end{grammarlens}

\subsection*{Your turn}
\begin{itemize}
\raggedright
  \item \emph{Say it:} \emph{gallus}
  \item \emph{Say it:} \emph{gallus}, once more
  \item \emph{Say it:} say \emph{gallus}
  \item \emph{Recall:} say the Latin for a dolphin, then the Latin for a whale, then say \emph{gallus} again
\end{itemize}

\begin{tcolorbox}[breakable,colback=teal!4,colframe=teal!35!black,title={Before you move on}]
What does gallus, gallī mean? (\textquotedblleft{}A cockerel\textquotedblright{}.) Say it once more.
\end{tcolorbox}
\section[labrum, labrī]{labrum, labrī — a lip}

\label{lesson:LA-C135-labrum}



\begin{quote}
Before the new one: say the Latin for a whale, then the Latin for a cockerel.
\end{quote}

\subsection*{You'll want to know: labrum, labrī}
\textbf{labrum, labrī} — \textquotedblleft{}a lip\textquotedblright{}.

\begin{grammarlens}[title={What you've built}]
One more word for this chapter.
\end{grammarlens}

\subsection*{Your turn}
\begin{itemize}
\raggedright
  \item \emph{Say it:} \emph{labrum}
  \item \emph{Say it:} \emph{labrum}, once more
  \item \emph{Say it:} say \emph{labrum}
  \item \emph{Recall:} say the Latin for a whale, then the Latin for a cockerel, then say \emph{labrum} again
\end{itemize}

\begin{tcolorbox}[breakable,colback=teal!4,colframe=teal!35!black,title={Before you move on}]
What does labrum, labrī mean? (\textquotedblleft{}A lip\textquotedblright{}.) Say it once more.
\end{tcolorbox}
\section[mentum, mentī]{mentum, mentī — the chin}

\label{lesson:LA-C135-mentum}



\begin{quote}
Before the new one: say the Latin for a cockerel, then the Latin for a lip.
\end{quote}

\subsection*{You'll want to know: mentum, mentī}
\textbf{mentum, mentī} — \textquotedblleft{}the chin\textquotedblright{}.

\begin{grammarlens}[title={What you've built}]
One more word for this chapter.
\end{grammarlens}

\subsection*{Your turn}
\begin{itemize}
\raggedright
  \item \emph{Say it:} \emph{mentum}
  \item \emph{Say it:} \emph{mentum}, once more
  \item \emph{Say it:} say \emph{mentum}
  \item \emph{Recall:} say the Latin for a cockerel, then the Latin for a lip, then say \emph{mentum} again
\end{itemize}

\begin{tcolorbox}[breakable,colback=teal!4,colframe=teal!35!black,title={Before you move on}]
What does mentum, mentī mean? (\textquotedblleft{}The chin\textquotedblright{}.) Say it once more.
\end{tcolorbox}
\section[dorsum, dorsī]{dorsum, dorsī — the back}

\label{lesson:LA-C135-dorsum}



\begin{quote}
Before the new one: say the Latin for a lip, then the Latin for the chin.
\end{quote}

\subsection*{You'll want to know: dorsum, dorsī}
\textbf{dorsum, dorsī} — \textquotedblleft{}the back\textquotedblright{}.

\begin{grammarlens}[title={What you've built}]
One more word for this chapter.
\end{grammarlens}

\subsection*{Your turn}
\begin{itemize}
\raggedright
  \item \emph{Say it:} \emph{dorsum}
  \item \emph{Say it:} \emph{dorsum}, once more
  \item \emph{Say it:} say \emph{dorsum}
  \item \emph{Recall:} say the Latin for a lip, then the Latin for the chin, then say \emph{dorsum} again
\end{itemize}

\begin{tcolorbox}[breakable,colback=teal!4,colframe=teal!35!black,title={Before you move on}]
What does dorsum, dorsī mean? (\textquotedblleft{}The back\textquotedblright{}.) Say it once more.
\end{tcolorbox}
